\chapter{Dialogue, commitments, and burdens of proof}
\label{sec:p18-dialogue-burden}

A demand for reasons has a procedural context. A speaker defending an
assertion, a committee comparing actions, and an investigator testing a
hypothesis need not have the same task. Before deciding that someone has
failed to prove a point, identify the point, the person responsible for
establishing it, the required support, and the assessment stage.

This chapter develops a manual reconstruction and assessment method from
three identified presentations. The stages and locution guide use
McBurney, Hitchcock, and Parsons's \emph{The Eightfold Way of Deliberation
Dialogue}, in the author preprint dated December 31, 2004
\cite{McBurneyHitchcockParsons2004DDF}. The classification of dialogue and
placement of persuasion within deliberation use Walton's 2010
\emph{Burden of Proof in Deliberation Dialogs}
\cite{Walton2010DeliberationBurden}. The three-burden distinction follows
Prakken and Sartor's 2009 \emph{A Logical Analysis of Burdens of Proof}
\cite{PrakkenSartor2009Burdens}. Their differences remain explicit.

The Eightfold Way supplies an outer deliberation framework, called DDF
in its source. It leaves the internal persuasion game and the transfer
of commitments from that game to the outer dialogue unspecified
(author preprint, pp.~24--25). The locution guide here is an operational
reconstruction for keeping a public record; its confirmation convention
is stated explicitly. The worked exchange supplies its own local evidence
policy. No general evidence evaluator, guarantee of agreement, or formal
implementation of the source's unspecified embedded game is claimed.

\section{Identify the dialogue and its issue}
\label{sec:p18-dialogue-types}

\begin{hbdefinition}{def:p18-dialogue-purpose}{Dialogue type, governing question, and episode}
A \emph{governing question} states what the participants seek to settle.
A \emph{dialogue type} specifies a characteristic starting problem,
the participants' roles or aims, and a collective purpose.
An \emph{episode} is a portion of an exchange whose question and
applicable rules have been identified. Identifying a type reconstructs
an activity; it does not establish the speakers' sincerity.

Walton's seven-type presentation makes the following distinctions
(2010 author text, pp.~2--4).
\begin{description}
\item[Persuasion.] A difference of opinion motivates attempts to bring
others to a position. The collective purpose is resolving or clarifying
the disputed issue.
\item[Inquiry.] A need for proof motivates finding and checking evidence
in order to establish or reject a hypothesis.
\item[Discovery.] A need for explanation motivates developing and
defending candidate hypotheses so that one can be selected for testing.
\item[Negotiation.] Conflicting interests motivate proposals and
concessions. Each participant seeks an advantageous result; collectively,
the parties seek a settlement they can live with.
\item[Information seeking.] Someone needs information that another
may provide; the collective purpose is exchanging information.
\item[Deliberation.] A practical question requires a choice.
Participants coordinate goals and possible actions to select an
appropriate available course.
\item[Eristic.] Personal conflict motivates verbal attack. Walton
describes its collective role in terms of exposing deeper conflict;
that description does not make abuse evidence for a disputed claim.
\end{description}
The earlier six-type classification reproduced in the Eightfold Way
does not list discovery separately. This chapter retains Walton's
identified seven-type variant.
\end{hbdefinition}

An action question can occur in persuasion as well as deliberation.
A participant defending the thesis that a specific action must be taken,
against another's challenge, can be in persuasion. Participants initially
asking what they should do and generating alternatives together can be
in deliberation. A fact--action distinction alone does not classify
the exchange.

\begin{hbmethod}{def:p18-dialogue-episode-method}{Reconstruct the procedural question}
Record the governing question, participants, current episode, and its
closing criterion. Identify assertions, suggested actions, and the
commitments that the context supports attributing to each speaker.
When an episode changes type, record its new question and its relation
to the outer purpose.

An \emph{embedded} inquiry or persuasion episode can support deliberation:
participants may need to establish whether an option satisfies a cost
constraint. Its result returns as an identified finding or concession.
A \emph{misleading shift} changes the task so as to evade an obligation,
as when a challenged proponent demands that the questioner prove the
opposite. Assess the shift against the agreed roles and purpose.
Changing type is not by itself an error.
\end{hbmethod}

Persuasion also distinguishes \emph{dispute}, in which parties defend
opposed theses, from \emph{dissent}, in which a proponent defends a thesis
and the respondent may merely question whether it has been established.
A respondent in dissent does not acquire the task of proving the
contrary merely by asking for reasons (Walton, 2010 author pp.~10--12).

\section{Commitments and the outer deliberation procedure}
\label{sec:p18-dialogue-commitments}

\begin{hbdefinition}{def:p18-dialogue-record}{Public commitments and typed contents}
A \emph{public commitment} is an attributable dialectical undertaking:
a participant has put a content on record and is responsible for it
under the applicable rules. It need not coincide with private belief.
Endorsing an action within a dialogue need not be a promise to execute
that action outside it.

DDF uses six mutually exclusive content types:
\emph{action}, \emph{goal}, \emph{constraint}, \emph{perspective},
\emph{fact}, and \emph{evaluation}.
An action is a candidate course of conduct; a goal is a desired outcome;
a constraint limits admissible choices; a perspective supplies a way
of assessing them. A fact is a proposition about a possible external
state: this type name does not certify truth. An evaluation assesses
an action in relation to a goal, constraint, or perspective.

A typed commitment is recorded as \((\tau,t)\), with content \(t\)
of type \(\tau\). A preference entry \((\mathrm{prefer},a,b)\) records
its speaker's preference for action \(a\) over action \(b\).
Each participant has a separate public store. Only that participant's
own moves alter it; others can read and challenge its entries.

Keep a chronological move record as well as current stores.
The record identifies the phase, active participants, recommendation,
and pending challenges. It also records prior proposals and assertions:
a current store alone cannot reveal whether a new utterance would
repeat an earlier one that has since been withdrawn.
\end{hbdefinition}

These conventions follow the Eightfold Way, author pp.~8--15 and
Appendix pp.~28--32. Its atomic language does not itself close stores
under logical consequence. Apart from replacing an action commitment
in the recommendation procedure, the guide does not enforce consistency
of everything a participant says. A preference can persist after a
supporting evaluation is retracted. Flag the unsupported retained
preference; deleting it on its owner's behalf would change the record.

\begin{hbmethod}{def:p18-dialogue-eight-stages}{Eight stages and their recurrence}
Organize a DDF deliberation through the following stages
(author pp.~5--8).
\begin{description}
\item[Open.] State the governing question and identify participants.
One agent opens; another's entry makes a dialogue between at least two.
\item[Inform.] Supply goals, constraints, perspectives, and facts.
\item[Propose.] Put candidate actions into consideration.
\item[Consider.] Evaluate and compare actions.
\item[Revise.] Change options or considerations in light of discussion.
\item[Recommend.] Move an action for the participants' responses.
\item[Confirm.] Obtain final endorsements of the recommendation.
\item[Close.] Record the result and end participation.
\end{description}
Open occurs once. Inform precedes first entry into the other substantive
stages. Propose precedes first Consider, Revise, Recommend, or Confirm;
Consider precedes first Revise; Confirm follows Recommend.
The intervening stages may recur, including renewed information gathering
and option development. Unanimous confirmation permits successful
closing. Withdrawal can instead end a dialogue without agreement.
If more than two participants remain, a departure need not end it;
if only two remain and one leaves, the withdrawal rule terminates it.

These are normative organization rules for the identified computational
framework, not an observed fixed chronology of human conversation.
They permit cycling. Neither finite termination nor objective optimality
of an agreed action follows from these stage requirements.
\end{hbmethod}

Walton's opening--argumentation--closing division is coarser than these
eight stages. His proposed account locates substantive burden-bearing
persuasion in argumentation, especially when an option becomes an
advocated recommendation. Early brainstorming can introduce a possibility
without undertaking to establish that it should be selected.
This is Walton's normative account, not a universal restriction on
when a forum may request reasons.

In particular, Walton argues that deliberation's opening need not
assign one side a global burden for a thesis: there may yet be no
fixed advocated thesis. Once someone advocates an option, objections
can require support in the relevant persuasive episode.
A comparative claim that an option is best remains vulnerable to
a better alternative or an overlooked harmful consequence.
Walton's discussion allows the required response to a safety objection
to depend on the seriousness and likelihood of harm and the competing
benefits (2010 author pp.~17--20). It does not supply a universal
numerical threshold or make a suggestion alone an obligation to
disprove every imaginable danger.

\subsection{The ten outer locutions}

The following guide preserves the substantive prerequisites and public
effects of Appendix L1--L10 of the held preprint, using the explicit
confirmation convention below. It describes recorded speech acts,
not a proof calculus. Subsequent speakers must have opened or entered
and still be participating. A move with no commitment effect remains
an event in the chronological record.

\begin{hbmethod}{def:p18-dialogue-locutions}{Locutions, prerequisites, and effects}
\begin{enumerate}
\item \texttt{open\_dialogue}: open the governing question, provided
this dialogue has not already been opened. No store entry is added.
\item \texttt{enter\_dialogue}: join the question already opened by
another participant. No store entry is added.
\item \texttt{propose}: offer a typed content that nobody has previously
proposed as that typed content. Before an action proposal, some
participant must have proposed or asserted a goal, constraint,
perspective, or fact. A proposal adds no commitment and requires
no immediate response.
\item \texttt{assert}: assert a typed content not previously asserted
by this speaker as an ordinary assertion. An evaluation requires
its action to have been proposed or asserted by some participant.
Add the typed entry to this speaker's store. An action assertion
answering another's recommendation replaces the speaker's other
action commitments. The separate confirmation role below permits
its recorded repeat endorsement.
\item \texttt{prefer}: prefer action \(a\) over action \(b\), provided
evaluations of both have already been asserted, possibly by others.
Add \((\mathrm{prefer},a,b)\) to this speaker's store; a different
preference is not thereby removed.
\item \texttt{ask\_justify}: ask another participant to justify a typed
content that participant asserted and has not retracted.
The question changes neither store. Its target must retract, defend
the assertion in embedded persuasion, or use embedded persuasion to
defend an exemption from justifying it. Announcing an exemption
does not establish one.
\item \texttt{move}: recommend an action already proposed or asserted
and still available on the record. Add it to this speaker's store
and remove that speaker's other action commitments. Every other
participant must respond by asserting or rejecting that action;
the source permits no third response to this recommendation.
Reasons for rejection can be pursued in a subsequent substantive episode.
\item \texttt{reject}: respond to another's recommendation by rejecting
its action. Remove that action from the rejecting speaker's store
if present. Do not edit the recommender's store. Rejection does not
assert a negative factual proposition.
\item \texttt{retract}: withdraw the commitment associated with one's
own unretracted earlier assertion, recommendation, or preference.
Remove the corresponding entry from one's own store. This rule does
not retract a proposal, which never inserted a commitment.
The historical proposal remains identifiable.
\item \texttt{withdraw\_dialogue}: leave a dialogue from which one has
not already withdrawn. Commitments are not erased. If two participants
remained immediately before this withdrawal, terminate the dialogue.
\end{enumerate}
\end{hbmethod}

The \emph{confirmation convention} here permits each participant one
confirmation endorsement of the current recommendation in its Confirm
round, even if they already asserted that action in response to the
recommendation. Record confirmation separately; adding an action already
present leaves the store unchanged. A revised recommendation requires
a new round. Ordinary repeated assertions outside this role remain
subject to the preprint's no-prior-assertion restriction. Retraction
does not erase an assertion's historical occurrence.
This stated operational convention makes the repeated endorsement
stage usable alongside the restriction on ordinary repeats.

\begin{proofboundary}{def:p18-dialogue-outer-boundary}{The embedded-procedure interface}
These locutions do not determine witness reliability, expert competence,
an evidential threshold, or the closing rule of embedded persuasion.
Nor do they automatically import every subdialogue commitment into
outer stores. Those matters require a selected embedded procedure.
Record its issue, roles, response obligations, evidence policy, and
closure convention. Import only the assertions, concessions, retractions,
or findings that the stated convention authorizes.
If the transcript lacks such a convention, identify the missing
information and limit the assessment accordingly.
\end{proofboundary}

\section{Three burdens and their different changes}
\label{sec:p18-dialogue-three-burdens}

\begin{hbdefinition}{def:p18-dialogue-three-burdens}{Persuasion, production, and tactical burdens}
Prakken and Sartor distinguish three burdens (2009 author text,
\S2, pp.~2--6, and \S4.1, pp.~18--19).

The \emph{burden of persuasion} assigns a party the risk that an issue
remains unestablished at the required standard when the proceeding
closes. Record the issue, owner, standard, and evaluation stage.
Failure to establish \(p\) to the standard does not alone establish
\(\neg p\).

The \emph{burden of production} requires enough evidential support
to make an issue eligible for consideration on its merits.
Where a minimally adequate evidential argument is required, bare
assertion is insufficient. Failure has the procedural consequence
specified by the forum; in the authors' legal model, the issue is
decided against the party whose production burden is unmet.
Passing this threshold permits consideration, without guaranteeing
final success.

A \emph{tactical burden} arises when a participant assesses that their
position would lose if the current record were the final record.
They have a strategic reason to add evidence or argument.
This is a stage-relative assessment, not a newly assigned legal burden.
What suffices depends on the role: establishing a proposition can
require more than raising enough doubt about it.
\end{hbdefinition}

\begin{hbmethod}{def:p18-dialogue-burden-ledger}{Keep separate issue-specific entries}
For each issue record:
\begin{enumerate}
\item the persuasion owner, required standard, and assessment on the
current record and at closure;
\item the production owner, threshold, whether it has been met, and
the consequence of failure;
\item the current tactical forecast, its assessor, and the response
that could affect the outcome.
\end{enumerate}
In Prakken and Sartor's modeled stages, production and persuasion
allocations remain fixed. A completed production step does not migrate
to the opponent whenever evidence changes. Tactical pressure can
move repeatedly.

A forum may authorize a separate reallocation decision, or participants
may open a different issue with another owner. Record the authority,
decision, and new allocation. A strong counterargument does not itself
reassign the original persuasion burden. The authors acknowledge
qualifications in actual legal practice; fixed allocation here is
not asserted as a universal rule for jurisdictions.
\end{hbmethod}

The allocation unit is an issue, not a permanent label on a person.
A respondent may challenge \(p\) without asserting \(\neg p\), then
assert \(r\) and become responsible for supporting \(r\).
These obligations coexist. Responsibility for \(r\) does not make
that respondent the original proponent of \(p\); questioning \(r\)
does not remove the original proponent's obligation.

Walton groups production and tactical burdens under a broad local-burden
heading (2010 author pp.~4--6). This chapter uses that heading only
for Walton's placement of obligations within episodes. When tracking
a change of burden, it retains Prakken and Sartor's sharper distinction
between evidential eligibility and the strategic need to respond.

\subsection{Challenges, concessions, and standards}

An unanswered challenge may leave an assertion unsupported or a
response obligation unmet. Its precise consequence depends on the
procedure. It does not automatically establish a negation, erase a
commitment, or transfer persuasion responsibility. DDF leaves stores
intact while requiring an allowed response. Without a response deadline
and default rule, it does not justify inventing a forfeiture.

Prakken and Sartor also discuss burdens of claiming, contesting, and
argument (\S5, author pp.~23--24). Claiming matters when a proposition
must be raised at a designated stage to count in one's favor.
Contesting matters where uncontested assertions count as concessions.
The burden of argument extends support requirements beyond factual
evidence to such claims as the rules used in deciding an issue.
Silence counts as a concession only under the applicable protocol.

The critical questions in
Chapter~\ref{sec:p18-argument-schemes-critical-questions}
identify relevant weaknesses and defeaters, while their procedural
effects still need specification. Asking about an expert's field can
request a missing premise. Alleging a particular undisclosed conflict
introduces a claim that itself needs evidence. Neither is adequately
described as every question transferring a single indivisible burden.
The sources' historical legal examples do not establish current law,
and named legal standards have no universal numerical translation.

\section{A complete deliberation with an embedded dissent}
\label{sec:p18-dialogue-worked}

\begin{hbexample}{def:p18-dialogue-library-exchange}{Temporary catalogue access}
Eden, Noor, and Lin form a fictional library committee choosing
temporary catalogue access for a one-week migration. All documents,
costs, and institutional conventions below are stipulated.
The alternatives are \(M\), an outsourced mirror, and \(D\), a staffed
manual search desk. The goal \(g\) is week-long search access.
The constraint \(c\) is a total cost at most \(1000\) credits.
The committee compares these two service options under that constraint
and the coverage requirement. It requires recommendation responses
followed by unanimous confirmation.

Use the following typed contents.
\begin{description}
\item[\(g,c\).] The goal and constraint just specified.
\item[\(q\).] The evaluation that \(M\) meets the total-cost limit.
\item[\(f_1\).] A supplier quotation prices its listed mirror services
at \(800\) credits and presents the list as complete.
\item[\(f_2\).] An applicable schedule requires an additional
\(300\)-credit integration fee omitted from that quotation.
\item[\(d\).] The desk meets the coverage requirement for a total
of \(450\) credits.
\item[\(h\).] The staff roster covers all stated service hours.
\item[\(m\).] With the fee included, the mirror costs
\(800+300=1100\) credits and exceeds the limit.
\end{description}
Here \(f_1,f_2,h\) have type fact; \(q,d,m\) have type evaluation.
\(f_1\) reports what the quotation lists and represents, without
guaranteeing that its list includes every applicable charge.
Learning \(f_2\) therefore need not require retracting \(f_1\).

The local dissent over \(q\) has this policy. Eden owns both the
persuasion and production burdens for \(q\).
Production requires an identifiable relevant supplier
cost record adequate to raise the issue. Final acceptance requires
an applicable itemized cost record supporting the limit with material
contrary cost evidence addressed. The committee provisionally accepts
an apparently complete quotation unless a material omission is evidenced.
The respondent may challenge \(q\) without undertaking to prove its
contrary. An affirmative counterclaim requires its own supporting record.
The episode closes when Eden withdraws \(q\), when its challengers
explicitly concede it after responses and no challenge remains pending,
or at an announced assessment after an opportunity to respond.
An unresolved or withdrawn \(q\) is recorded as unestablished;
\(\neg q\) is not automatically recorded.

The embedded episode admits a cost document, a question about its
applicability or completeness, a supported counterclaim, a response
or concession, and withdrawal. Documents attach to assertions;
they are not an eleventh DDF locution. In this example, speakers'
assertions during the episode are explicitly entered in their outer
stores. A finding is not silently attributed to everyone as a personal
assertion. This local policy is an added manual convention for the
example, not the unspecified persuasion game of the source.
\end{hbexample}

\subsection{The transcript and every public effect}

Write \(+x\) and \(-x\) for adding and removing the typed content \(x\)
from the named speaker's store. The types were fixed above;
\(+D\) adds the action entry. Write \(+\pi\) for adding that speaker's
preference for \(D\) over \(M\). Zero change still records the event.
One unsuccessful attempted move is explicitly identified.

\begin{enumerate}
\item \emph{Open:} Eden opens the governing question. Store change:
zero. No specific option is yet a thesis Eden must establish.
\item Noor enters. Store change: zero.
\item Lin enters. Store change: zero. All three participate.
\item \emph{Inform:} Eden proposes \(g\). Store change: zero.
\item Eden asserts \(g\). Eden: \(+g\).
\item Noor asserts \(g\). Noor: \(+g\).
\item Lin asserts \(g\). Lin: \(+g\).
\item Eden asserts \(c\). Eden: \(+c\).
\item Noor asserts \(c\). Noor: \(+c\).
\item Lin asserts \(c\). Lin: \(+c\). Goal and constraint are
now commitments of each participant.
\item \emph{Propose:} Eden proposes \(M\). Store change: zero.
The prior goal and constraint satisfy the proposal prerequisite.
\item Noor proposes \(D\). Store change: zero.
\item \emph{Consider:} Eden asserts \(q\): \emph{The mirror fits our
budget.} Eden: \(+q\). The prior proposal of \(M\) permits its evaluation.
An advocated assessment has now replaced mere suggestion.
\item Noor asks Eden to justify \(q\): \emph{Does that include every
mandatory charge?} Store change: zero. The target is unretracted.
Noor has not asserted \(\neg q\).
\item Eden attempts: \emph{You prove that it does not fit.}
Store change: zero. This neither retracts nor defends \(q\), nor
justifies an exemption. It cannot be a challenge to Noor's \(\neg q\),
because Noor never asserted that content. The original challenge
remains pending.
\item \emph{Embedded dissent:} the participants record \(q\), Eden's
role, and the stated local policy. Store change: zero.
The issue is whether Eden's assessment has been established.
\item Eden asserts \(f_1\) and attaches the quotation.
Eden: \(+f_1\). This meets the production threshold.
If the record closed with this apparently complete quotation and
no contrary cost evidence, the stipulated policy would accept \(q\).
Noor therefore has a tactical reason to object if seeking its rejection.
Eden retains persuasion responsibility.
\item Noor asserts \(f_2\) and supplies the applicable schedule,
including the provision making its fee mandatory. Noor: \(+f_2\).
This is a supported counterclaim. The record no longer supports
accepting \(q\). Tactical pressure returns to Eden to address the
omission or withdraw. The already-met production burden has not
been reassigned to Noor.
\item Eden retracts \(q\): \emph{I withdraw the claim that the mirror
meets the limit.} Eden: \(-q\).
The report \(f_1\) remains accurate as a report of the quotation.
The dissent closes with \(q\) withdrawn and unestablished.
The challenge has an allowed response. Nobody automatically acquires
\(\neg q\) in their store.
\item \emph{Return to Consider:} Noor asserts \(d\) with the
\(450\)-credit staffing calculation and roster. Noor: \(+d\).
This evaluates the previously proposed \(D\).
\item Lin checks the roster and asserts \(h\). Lin: \(+h\).
Its stipulated coverage supports the coverage part of \(d\);
a bare unsupported assertion of \(h\) would not verify it.
\item Eden asserts \(m\), explains \(800+300=1100>1000\), and
identifies the quotation and mandatory schedule as its premises.
Eden: \(+m\). This new evaluation does not repeat the retracted \(q\).
\item Noor prefers \(D\) over \(M\). Noor: \(+\pi\).
The evaluations \(d,m\) satisfy the preference prerequisite.
The stated basis is coverage and cost.
\item Eden records that preference. Eden: \(+\pi\).
\item Lin records that preference. Lin: \(+\pi\).
Each is a separate attributable commitment.
\item \emph{Recommend:} Noor moves \(D\). Noor: \(+D\), removing
any other action commitment; none is present. Eden and Lin owe
the specified recommendation response.
\item Eden asserts \(D\) in response. Eden: \(+D\), replacing
any other action commitment; none is present.
\item Lin asserts \(D\) in response. Lin: \(+D\), with the same
replacement convention. Required recommendation responses are complete.
\item \emph{Confirm:} Noor confirms \(D\). The store already contains
\(D\), so its contents do not change. Record the confirmation.
\item Eden confirms \(D\). Store change: zero. This is the explicitly
permitted phase-specific repeat endorsement.
\item Lin confirms \(D\). Store change: zero. With all three
confirmations recorded, the committee records \(D\) as its unanimous
decision and enters Close.
\item \emph{Close:} Eden withdraws. Commitments remain.
Noor and Lin are the two remaining participants.
\item Noor withdraws. Commitments remain. Only Lin would remain,
so the withdrawal condition terminates the session. The recorded
decision is retained.
\end{enumerate}

Suppressing type tags for readability, Eden's final store is
\(g,c,f_1,m,\pi,D\); Noor's is \(g,c,f_2,d,\pi,D\);
Lin's is \(g,c,h,\pi,D\). The retracted \(q\) appears in history,
but not Eden's current store. Documents and the collective decision
remain distinct records. None of these stores inventories private belief.

\subsection{The resulting assessment and two failure branches}

Proposing \(M\) carried no commitment. Advocating its cost assessment
did. The quotation raised an issue without making its conclusion immune
to later evidence. The supported fee changes the cost calculation:
\[
800+300=1100>1000,\qquad 450\leq1000.
\]
The roster provides the separate coverage premise.
The comparison supports the desk under the stipulated records and
criteria. It establishes neither that every possible alternative
was considered nor optimality under unstated values.

Eden's persuasion responsibility for \(q\) remains fixed.
The production threshold is passed at line 17. The forecast of losing
without a response changes at lines 17 and 18. Noor separately takes
responsibility for \(f_2\) and supports it. These different obligations
do not turn Noor's original question into a duty to prove the contrary
of every aspect of Eden's assessment.

\begin{hbexample}{def:p18-dialogue-unanswered-challenge}{If the proponent never answers}
After line 14, suppose Eden keeps demanding proof of the contrary,
without a quotation, defense, exemption argument, or retraction.
Eden still publicly holds \(q\). A current store therefore does not
by itself show adequate support. The response remains pending and
the stipulated production threshold remains unmet.

At an announced closing after an opportunity to respond, the local
policy records \(q\) as unestablished. It does not infer that the
mirror necessarily exceeds the limit. Without a supplied closing
or deadline rule, the outer guide leaves an unfinished dialogue,
not an automatic victory. Noor's silence in response to the
misdirected demand neither concedes \(q\) nor defends it.
\end{hbexample}

\begin{hbexample}{def:p18-dialogue-rejection-branch}{Rejection and incomplete consensus}
Instead, suppose Lin rejects \(D\) at line 28. This removes \(D\)
from Lin's store if present and leaves Noor's and Eden's stores alone.
It does not assert that the desk fails, the mirror succeeds, or a
document is false. Pursue Lin's reasons in renewed Consider or Revise.
Two endorsements do not satisfy this framework's unanimous decision
procedure.

A majority-vote forum could choose another closing rule, but that is
a different stated variant. A stored preference also does not replace
the required recommendation response and confirmation.
\end{hbexample}

\section{Assessment limits and further source developments}

The method produces an episode map, chronological move record, separate
commitment stores, issue-specific burden ledger, evidence and challenges,
and a qualified closing assessment. These outputs expose an evaded
response, an attributed assertion never made, or a premature declaration
of agreement.

The sources have different roles. The Eightfold Way proposes a framework
for computational deliberation. Walton analyzes the normative placement
of burdens. Prakken and Sartor distinguish burdens through stages of
defeasible proof. The fictional committee adds a manual evidence policy.
None of these roles establishes that people follow these rules or that
the rules improve actual decisions.

\begin{proofboundary}{def:p18-dialogue-next-boundary}{Separate source developments}
The supplied extent is the identified dialogue distinctions, complete
ten-locution outer guide, three-burden method, and worked exchange with
its explicit embedded policy. Prakken and Sartor's full LIS argument
game, including issue-dependent roles and its underlying defeasible
calculus, remains a separate development. So do a fully specified
persuasion game filling DDF's interface, formal strategy or termination
results, negotiation and inquiry protocols, and competing accounts
of presumptions, defaults, and burden allocation.

The source's normative desiderata are not converted into unproved
formal guarantees. Institutional applications need their own identified
procedure and sources. Existing abstract and structured argumentation
semantics can assess explicitly constructed argument graphs; merely
assigning a graph to utterances does not recover commitment updates,
burden allocation, or closing rules.
\end{proofboundary}
